%% ====================================================================
\section{Microbenchmarks}
\label{sec:microbenchmarks}
%% ====================================================================

This section covers per-operation sign and verify costs for every signing primitive in the Lux production stack. Numbers are from \texttt{github.com/luxfi/consensus} \cite{consensusgo}, \texttt{github.com/luxfi/crypto/$\langle$bls,mldsa,slhdsa$\rangle$}, \texttt{github.com/luxfi/pulsar} \cite{ringtailgo}, and \texttt{github.com/luxfi/lens} \cite{lensgo}, on Apple M1 Max single core.

\subsection{Single-signature primitives}

\begin{tabular}{@{}lrrr@{}}
  \toprule
  \textbf{Scheme} & \textbf{Sign} & \textbf{Verify} & \textbf{Verify (cached)} \\
  \midrule
  secp256k1 (C, native)            & 658 ns        & 658 ns        & --- \\
  Ed25519 (Go stdlib)              & ---           & 205 \textmu s & 1.1 \textmu s \\
  P-256 (Go stdlib)                & ---           & 121 \textmu s & 1.0 \textmu s \\
  ML-DSA-44                        & ---           & 140 \textmu s & --- \\
  ML-DSA-65                        & 495 \textmu s & 250 \textmu s & 3.0 \textmu s \\
  ML-DSA-65 (via UTXO Fx context)  & ---           & 254 \textmu s & 3.0 \textmu s \\
  ML-DSA-87                        & ---           & 420 \textmu s & --- \\
  SLH-DSA-SHA2-192f                & ---           & 1.92 ms       & 131 \textmu s \\
  BLS12-381 (single)               & 350 \textmu s & 820 \textmu s & --- \\
  \bottomrule
\end{tabular}

Reproduction: \texttt{cd consensus \&\& GOWORK=off go test -bench=. -benchtime=1s ./crypto/bls/...} (and analogous for \texttt{mldsa}, \texttt{slhdsa}). Cached verify uses a process-local LRU; the cache hit is dominated by the LRU lookup, not the cryptographic verify.

\subsection{BLS aggregate verify}

\begin{tabular}{@{}lrr@{}}
  \toprule
  \textbf{Operation} & \textbf{Time} & \textbf{Source} \\
  \midrule
  BLS aggregate (100 sigs) & 5.3 ms & \texttt{protocol/quasar BenchmarkBLSAggregation/100} \\
  BLS aggregated verify (100 signers) & 875 \textmu s & \texttt{protocol/quasar BenchmarkBLSAggregatedVerification/100} \\
  \bottomrule
\end{tabular}

The aggregated verify is constant in signer count (the aggregate signature is 96 bytes regardless of $n$); the cost is dominated by the pairing computation. Reproduction: \texttt{GOWORK=off go test -bench=BLSAggregat -benchtime=2s ./protocol/quasar/...}.

\subsection{Pulsar (lattice threshold) per-party costs}

Measured timings on a single Apple M1 Max P-core (Darwin 25.4.0
arm64, Go 1.26.2), Go canonical, via \texttt{cd pulsar \&\& go run . l 5
K} which calls \texttt{sign.LocalRun(5)} (\texttt{pulsar/sign/local.go:21}).
Per-phase median across 5 end-to-end runs at $K = 5$ and $K = 21$
(2026-05-04):

\begin{center}
\begin{tabular}{@{}lrr@{}}
  \toprule
  \textbf{Phase} & \textbf{$K = 5$ median (ms)} & \textbf{$K = 21$ median (ms)} \\
  \midrule
  \texttt{Gen} (one party)                  & 3.11    & 24.49  \\
  \texttt{SignRound1} (per party)           & 17.40   & 171.19 \\
  \texttt{SignRound2Preprocess} (per party) & 67.30   & 581.93 \\
  \texttt{SignRound2}  (per party)          & 2.52    & 14.65  \\
  \texttt{SignFinalize}                     & 0.065   & 0.148  \\
  \texttt{Verify}                           & 0.94    & 1.49   \\
  \texttt{Total Signing} (per party)        & 87.37   & 763.58 \\
  \bottomrule
\end{tabular}
\end{center}

Round 2 Preprocess dominates the per-party cost --- it scales roughly
linearly in $K$ (the MAC-verification matrix multiply is $O(K^2)$ in
the matrix size but dominated here by per-party setup). Round 1 is
offline-preprocessable: it is independent of the message $m$, so a
node can run it ahead of time on quiescent CPU and store the output
until message-time. The C++ port at LP-137 \cite{lp137} closes the bulk
of the matrix-multiply gap; GPU acceleration is targeted to bring
Round 2 Preprocess below 100 ms at $K = 21$ (Section~\ref{sec:gpu-headroom}).
Reproduction: \texttt{cd pulsar \&\& GOWORK=off go run . l 5 21}.

\paragraph{Pulsar NTT C++ microbench (M1 Max).}
\texttt{luxcpp/crypto/build/pulsar/pulsar\_sign\_ntt\_bench} (single
NTT over $\mathbb{Z}/Q\mathbb{Z}$ with $Q = 0x1000000004A01$, $N = 256$,
LP-073 canonical Pulsar ring) reports the following CPU per-NTT cost
(2026-05-04, M1 Max, batches independent):

\begin{center}
\begin{tabular}{@{}lrrr@{}}
  \toprule
  \textbf{Batch} & \textbf{Reps} & \textbf{Median (\textmu s, total)} & \textbf{Per-NTT (\textmu s)} \\
  \midrule
  1   & 2000 & 1.37   & 1.375 \\
  8   & 500  & 10.71  & 1.339 \\
  32  & 200  & 43.04  & 1.345 \\
  128 & 50   & 172.00 & 1.344 \\
  \bottomrule
\end{tabular}
\end{center}

The flat per-NTT $\sim 1.34$ \textmu s scaling shows the SIMD path is
saturated; further wins come from GPU dispatch (\texttt{LUX\_PULSAR\_NTT\_BACKEND=metal}
or \texttt{=cuda}; both unavailable on this host). Reproduction:
\texttt{./build/pulsar/pulsar\_sign\_ntt\_bench} from
\texttt{luxcpp/crypto/build/}.

\subsection{Lens (curve threshold) per-party costs}

Lens implements the FROST 2-round Schnorr protocol \cite{rfc9591} over three groups: Ed25519, secp256k1, and Ristretto255. Indicative single-party signing on M1 Max:

\begin{tabular}{@{}lrrr@{}}
  \toprule
  \textbf{Group} & \textbf{Sign1 (per party)} & \textbf{Sign2 (per party)} & \textbf{Aggregate + Verify} \\
  \midrule
  Ed25519        & 90 \textmu s   & 110 \textmu s  & 320 \textmu s \\
  secp256k1      & 75 \textmu s   & 95 \textmu s   & 280 \textmu s \\
  Ristretto255   & 95 \textmu s   & 115 \textmu s  & 340 \textmu s \\
  \bottomrule
\end{tabular}

Lens is roughly $\sim 1000\times$ faster than Pulsar on a per-party basis (curve scalar mul vs lattice matrix-vector multiply). The trade-off is: Lens is classical (curve discrete log), broken in polynomial time by Shor; Pulsar is post-quantum (Module-LWE). Lens is the right primitive for control-plane operations (governance, parameter updates) where compactness and speed dominate; Pulsar is the right primitive for the consensus PQ floor. Reproduction: \texttt{cd lens \&\& go test -bench=. ./sign/...}.

\subsection{Hash primitives}

The Pulsar-SHA3 hash suite uses FIPS-202 / SP 800-185 primitives:

\begin{tabular}{@{}lr@{}}
  \toprule
  \textbf{Primitive} & \textbf{Throughput (single core)} \\
  \midrule
  SHA3-256 (\texttt{golang.org/x/crypto/sha3}) & 250 MB/s \\
  cSHAKE256 (32-byte output)                   & 220 MB/s \\
  TupleHash256                                 & 200 MB/s \\
  KMAC256                                      & 240 MB/s \\
  BLAKE3 (1 MB)                                & 4.2 GB/s \\
  \bottomrule
\end{tabular}

The Pulsar-SHA3 suite is the production default; Pulsar-BLAKE3 is faster but non-normative. Reproduction: \texttt{go test -bench=. ./hash/...} from the pulsar root.

\subsection{Z-Chain Groth16 wrapper}

The optional Groth16 wrapper around the ML-DSA cert set (LP-307 \cite{lp307}) compresses verification of $N$ ML-DSA-65 signatures into a $\sim 192$-byte SNARK plus a 32-byte public-inputs hash:

\begin{tabular}{@{}lr@{}}
  \toprule
  \textbf{Operation} & \textbf{Time} \\
  \midrule
  Groth16 prover (ML-DSA-65 verification circuit, $\sim 2^{22.5}$ R1CS constraints) & $\sim 400$ ms CPU / $\sim 5$--$15$ ms GPU \\
  Groth16 verifier (single proof) & 1--3 ms (pairing-dominated) \\
  Groth16 verifier (batch of 64 via accel) & $\sim 200$--$500$ \textmu s/proof \\
  \bottomrule
\end{tabular}

The Groth16 wrapper is classical (Remark~\ref{rem:groth16-wrapper-warp} of \cite{luxwarpv2}): pairing-based, broken by Shor. PQ finality lives in ML-DSA + Pulsar, not in the Groth16 proof. The wrapper is useful for compression and EVM verifier compatibility only.

\subsection{Note on the stale 357 \textmu s claim}

Earlier Lux drafts (lux-triple-proof-consensus, lux-master-security-model, lux-performance-security-tradeoffs) carried a ``357 \textmu s epoch finality'' number. That number does not match any measured operation in the current code. Closest real candidates: BLS single keygen (350 \textmu s), ML-DSA-65 sign (495 \textmu s); Groth16 prover time is not 357 \textmu s in any reasonable cost model.

\textbf{Recommendation.} Cite the measured numbers in this report. The QuasarCert verify cost is $\sim 2$--$5$ ms CPU on M1 Max (Section~\ref{sec:prism-composition}). The 357 \textmu s figure is removed from the production benchmark surface.
